\documentclass[11pt,a4paper]{article}
\usepackage[margin=1in]{geometry}
\usepackage[T1]{fontenc}
\usepackage{lmodern}
\usepackage[hidelinks]{hyperref}
\pagestyle{empty}
\setlength{\parskip}{0.55em}
\setlength{\parindent}{0pt}

\begin{document}

\begin{center}
{\Large\bfseries Cover Letter}\\[0.4em]
\textit{Universe} (MDPI)\\[0.6em]
27 July 2026
\end{center}

\vspace{0.8em}

Dear Editors,

Please find enclosed our manuscript entitled
\textbf{``JW-FD: A Long Horizon Multimodal Solar Flare Forecasting Dataset''},
by Shao Mingfu, Lin Jiaben (corresponding author), Wang Hui, Tong Liyue, Yang Chen, Zhang Yin, and Li Yuyang, for consideration as an Article in the Special Issue ``New Discoveries in Astronomical Data (II)'' (Section: Astroinformatics and Astrostatistics), \textit{Universe}.

Solar flare forecasting is central to both fundamental heliophysics and operational space weather. Progress in data-driven methods has been limited by fragmented public resources that often cover a single solar cycle, release only images or only parameters, or define labels and splits inconsistently. JW-FD addresses this gap as shared infrastructure: a long-horizon, active-region-centered multimodal corpus that links magnetogram imagery, magnetic features, flare labels, and evolution videos under a leakage-aware evaluation protocol, with an open construction pipeline for reuse and extension. Its value lies less in any single model result than in enabling comparable, reproducible studies across statistics, classical machine learning, deep learning, and multimodal approaches spanning Cycles~24--25.

The manuscript presents a curated, machine-learning-ready solar dataset derived from long-term SDO/HMI and NOAA/SWPC observations, with standardized preprocessing, multimodal alignment, and evaluation splits intended to support reproducible data-driven heliophysics. This contribution aligns with the special issue emphasis on astronomical data processing, astroinformatics, and machine-learning methods that enable new scientific use of large observational archives, including time-domain solar data.

This manuscript has not been previously submitted to any MDPI journal.

We confirm that neither the manuscript nor any parts of its content are currently under consideration for publication with or published in another journal.
All authors have approved the manuscript and agree with its submission to \textit{Universe}.

Thank you for considering our submission.

\vspace{1.2em}
Sincerely,

\vspace{1.5em}
Lin Jiaben\\
Corresponding Author\\
State Key Laboratory of Solar Activity and Space Weather, NAOC\\
Email: \href{mailto:jiabenlin@bao.ac.cn}{jiabenlin@bao.ac.cn}

\end{document}
